% Generated draft: requires mathematical and source review.
\documentclass[11pt]{article}
\usepackage[margin=25mm]{geometry}
\usepackage{fontspec}
\setmainfont{Latin Modern Roman}
\usepackage{amsmath,amssymb,mathtools,unicode-math}
\setmathfont{Latin Modern Math}
\usepackage{enumitem,needspace,xurl,hyperref}
\setlength{\parindent}{0pt}
\setlength{\parskip}{5pt}
\setlength{\emergencystretch}{3em}
\newcommand{\N}{\mathbb{N}}
\newcommand{\Z}{\mathbb{Z}}
\newcommand{\Q}{\mathbb{Q}}
\newcommand{\R}{\mathbb{R}}
\newcommand{\C}{\mathbb{C}}
\newcommand{\F}{\mathbb{F}}
\newcommand{\A}{\mathbb{A}}
\newcommand{\PP}{\mathbb{P}}
\newcommand{\E}{\mathbb{E}}
\newcommand{\eps}{\varepsilon}
\newcommand{\dd}{\,\mathrm{d}}
\begin{document}
\section*{UnsolvedMath ID: 30004451}
\textbf{Source problem code:} catalogue entry
\subsection{Definitions and mathematical statement}
Let \(A\) be a separable nonunital complex \(C^*\)-algebra, and let \(M(A)\) be its multiplier algebra, with \(A\) identified with its canonical ideal in \(M(A)\). Define the corona algebra and quotient map by
\[
Q(A):=M(A)/A,\qquad \pi_A:M(A)\longrightarrow Q(A).
\]
Write \(\operatorname{Aut}(Q(A))\) for the group of complex-linear \(*\)-automorphisms of \(Q(A)\).

The supplied record does not define topological triviality. For this draft, use the standard Borel-graph convention, whose agreement with the intended source definition requires review. Equip \(M(A)\) with the strict topology generated by the seminorms
\[
p_a(m)=\|ma\|+\|am\|\qquad(a\in A),
\]
and let \(M(A)_1=\{m\in M(A):\|m\|\leq 1\}\) have the induced topology. For \(\alpha\in\operatorname{Aut}(Q(A))\), put
\[
\Gamma_\alpha=\{(m,n)\in M(A)_1\times M(A)_1:
\alpha(\pi_A(m))=\pi_A(n)\}.
\]
Call \(\alpha\) topologically trivial if \(\Gamma_\alpha\) is Borel for the product of the strict topologies. Set
\[
\mathcal T(A)=\{\alpha\in\operatorname{Aut}(Q(A)):
\alpha\text{ is topologically trivial}\}.
\]
Here CH is the continuum hypothesis \(2^{\aleph_0}=\aleph_1\), and PFA is the Proper Forcing Axiom: for every proper forcing partial order and every family of at most \(\aleph_1\) dense subsets, there is a filter meeting every member of the family.

\textbf{Question.} With the preceding interpretation of topological triviality, do both of the following assertions hold?
\[
\mathrm{CH}\ \Longrightarrow\
\left[\text{for every separable nonunital complex }C^*\text{-algebra }A,
\quad
\left|\operatorname{Aut}(Q(A))\setminus\mathcal T(A)\right|
=2^{\aleph_1}\right],
\]
\[
\mathrm{PFA}\ \Longrightarrow\
\left[\text{for every separable nonunital complex }C^*\text{-algebra }A,
\quad
\operatorname{Aut}(Q(A))=\mathcal T(A)\right]?
\]
The two implications concern separate set-theoretic assumptions. The universal quantifier over \(A\) is retained without adding approximation, stability, simplicity, or approximate-identity hypotheses.

\textbf{Supplied qualifications.} The local record labels the problem partially solved. Its supplied assessment says that PFA-side rigidity is proved for broad classes, including algebras with the metric approximation property and suitable approximate identities, and that OCA together with MA gives topological triviality. It also says that the full assertion for all separable nonunital algebras, especially the uniform CH cardinality assertion, remains unsettled. These are supplied dated curation claims, not independently verified conclusions. The record gives neither the precise class hypotheses nor the versions of OCA and MA needed to turn its background assertions into precise supplementary theorems.
\subsection{Short English statement}
For every separable nonunital complex \(C^*\)-algebra \(A\), ask whether CH forces its corona \(M(A)/A\) to have \(2^{\aleph_1}\) automorphisms that are not topologically trivial, while PFA forces every corona automorphism to be topologically trivial. This draft interprets topological triviality through a Borel graph on multiplier unit balls with their strict topologies; the intended source definition requires confirmation. The partially solved label and background results are supplied dated assessments.
\subsection{Sources}
\begin{itemize}
\item Local UnsolvedMath problem registry, permanent ID \(30004451\), title: \emph{Set-Theoretic Rigidity of Corona-Algebra Automorphisms}; collection: \emph{Oberwolfach Reports Open Problems}; exact source problem code: \texttt{catalogue entry}. The record gives proposed year 2020 and no named proposer. Its \texttt{corrected\_verified} statement label is accompanied by a qualification that only formula delimiters, notation, grammar, and typography were repaired, with mathematical content and scope preserved; it does not constitute independent mathematical verification here.
\item Supplied source citation: \emph{Set Theory (2021)}, \emph{Oberwolfach Reports}, DOI \texttt{10.4171/owr/2020/14}, pp.\ 797--855. \url{https://doi.org/10.4171/owr/2020/14}. The background describes these reports as records of mathematical workshops and their open-problem discussions. The same DOI is also listed as \emph{Oberwolfach Report 14/2020}; the differing year labels are retained as supplied.
\item Supplied literature reference: Coskey--Farah (2014). \url{https://doi.org/10.1090/S0002-9947-2014-06146-1}.
\item Supplied literature reference: McKenney--Vignati (2021). \url{https://arxiv.org/abs/1806.09676}.
\item Supplied literature reference: \emph{Oberwolfach Report 14/2020}. \url{https://doi.org/10.4171/owr/2020/14}.
\item Literature assessment date: \texttt{2026-08-22}. This is supplied dated curation, not independently verified. The record describes it as literature triage derived from the supplied Oberwolfach curation and requests independent verification of time-sensitive status and cited claims. Its assessment reports broad-class PFA rigidity results and an OCA-plus-MA topological-triviality result, while describing the full all-separable-algebras conjecture, especially its uniform CH cardinality assertion, as unsettled.
\end{itemize}
\end{document}
